\documentclass[11pt]{amsart}

\usepackage[margin=1in]{geometry}
\usepackage{amsmath,amssymb,amsthm,mathtools}
\usepackage[shortlabels]{enumitem}
\usepackage{hyperref}

\numberwithin{equation}{section}

\newtheorem{theorem}{Theorem}[section]
\newtheorem{proposition}[theorem]{Proposition}
\newtheorem{lemma}[theorem]{Lemma}
\newtheorem{corollary}[theorem]{Corollary}
\newtheorem{assumption}[theorem]{Assumption}
\theoremstyle{definition}
\newtheorem{definition}[theorem]{Definition}
\theoremstyle{remark}
\newtheorem{remark}[theorem]{Remark}

\newcommand{\E}{\mathbb E}
\newcommand{\Prob}{\mathbb P}
\newcommand{\N}{\mathbb N}
\newcommand{\R}{\mathbb R}
\newcommand{\Sph}{\mathbb S}
\newcommand{\Span}{\operatorname{span}}
\newcommand{\supp}{\operatorname{supp}}
\newcommand{\sgn}{\operatorname{sgn}}
\newcommand{\vc}{\operatorname{vc}}
\newcommand{\eps}{\varepsilon}
\newcommand{\abs}[1]{\lvert #1 \rvert}
\newcommand{\norm}[1]{\lVert #1 \rVert}
\newcommand{\ip}[2]{\langle #1,#2\rangle}
\newcommand{\distpm}[2]{d_{\pm}(#1,#2)}

\title{A Check of the VC Draft for the Probabilistic Modulus Gap}
\author{}

\begin{document}
\maketitle

\section{Setup}

Let \(\eta,\xi_2,\xi_3,\dots\) be independent real-valued random variables.
The exceptional variable \(\eta\) is centered and \(L^2\)-normalized:
\[
        \E\eta=0,\qquad \norm{\eta}_{L^2}=1.
\]
For a fixed \(\delta>0\), the good coordinates are assumed to satisfy
\[
        \E\xi_i=0,\qquad \norm{\xi_i}_{L^2}=1,\qquad
        \delta\le \norm{\xi_i}_{L^1}\le 1-\delta
        \qquad (i\ge2).
\]
For coefficient vectors \(a,b\in\ell^2(\N)\), write
\[
X_a=a_1\eta+\sum_{i\ge2}a_i\xi_i,\qquad
X_b=b_1\eta+\sum_{i\ge2}b_i\xi_i.
\]
In finite dimension we write
\[
Z=(\eta,\xi_2,\ldots,\xi_n)\in\R^n,\qquad
\langle Z,a\rangle=a_1\eta+\sum_{i=2}^n a_i\xi_i .
\]
Finally,
\[
        \distpm{u}{v}:=\min\{\norm{u-v}_2,\norm{u+v}_2\}.
\]

The question checked in this note is whether the one-exception hypotheses
above imply a uniform fixed-probability lower bound of the form
\[
\Prob\left(\bigl||X_a|-|X_b|\bigr|\ge\theta_0\right)\ge p_0
\]
for all orthogonal unit coefficient vectors \(a,b\).  Such a bound would be
exactly the small-ball input needed for the standard VC discretization
argument.

\section{VC Background}

The following standard VC facts are the only empirical-process input used in
the conditional argument below.

\begin{definition}[Shattering and VC dimension]
Let \(\mathcal F\) be a class of Boolean functions on \(\R^n\).  A finite set
\(S\subset\R^n\) is shattered by \(\mathcal F\) if for every map
\(g:S\to\{0,1\}\), there is \(f\in\mathcal F\) such that \(f|_S=g\).  The
VC-dimension \(\vc(\mathcal F)\) is the largest cardinality of a finite set
shattered by \(\mathcal F\).
\end{definition}

\begin{proposition}[VC law of large numbers]\label{prop:vc-lln}
Let \(\mathcal F\) be a class of Boolean functions with finite VC-dimension and
let \(X_1,\ldots,X_m\) be i.i.d. copies of a random vector \(X\).  There is an
absolute constant \(C>0\) such that
\[
\E\sup_{f\in\mathcal F}
\abs{\frac1m\sum_{j=1}^m f(X_j)-\E f(X)}
\le C\sqrt{\frac{\vc(\mathcal F)}{m}}.
\]
Consequently, for every \(t>0\), with probability at least \(1-e^{-t}\),
\[
\sup_{f\in\mathcal F}
\abs{\frac1m\sum_{j=1}^m f(X_j)-\E f(X)}
\le C\sqrt{\frac{\vc(\mathcal F)}{m}}+\sqrt{\frac{t}{2m}}.
\]
\end{proposition}

\begin{proposition}[Intersections and unions]\label{prop:vc-intersection}
Let \(\mathcal F_1,\ldots,\mathcal F_k\) be Boolean classes with
\(\vc(\mathcal F_i)=v_i\).  If \(\mathcal F\) is the class of indicators of
intersections of one set from each \(\mathcal F_i\), or the class of indicators
of unions of one set from each \(\mathcal F_i\), then
\[
\vc(\mathcal F)\le C\Bigl(\sum_{i=1}^k v_i\Bigr)\log(Ck),
\]
where \(C>0\) is absolute.
\end{proposition}

\begin{lemma}[Polynomial threshold classes]\label{lem:poly-vc}
Fix integers \(D,N_x,N_\theta\).  Let \(P(x,\theta)\) range through a family of
real polynomials of degree at most \(D\) in the variables
\((x,\theta)\in\R^{N_x}\times\R^{N_\theta}\), with arbitrary parameter
\(\theta\).  Then the Boolean class
\[
        \bigl\{\mathbf 1_{\{x:P(x,\theta)\ge0\}}:\theta\in\R^{N_\theta}\bigr\}
\]
has VC-dimension \(O_D(N_x+N_\theta)\).  By
Proposition~\ref{prop:vc-intersection}, the same remains true for
intersections of a bounded number of such classes.
\end{lemma}

\section{The Proposed Probabilistic Gap is False}

The attempted replacement for the Paley--Zygmund step was the following
assertion: there should exist constants
\(\theta_0=\theta_0(\delta,\eta)>0\) and
\(p_0=p_0(\delta,\eta)>0\) such that, for every orthogonal pair
\(a,b\in\ell^2(\N)\) with \(\norm{a}_2=\norm{b}_2=1\),
\[
\Prob\left(\bigl||X_a|-|X_b|\bigr|\ge\theta_0\right)\ge p_0 .
\]
This assertion does not follow from the hypotheses in Section~1.  In fact it
is false, even for coordinate vectors.

\begin{proposition}[Rare-spike obstruction]\label{prop:rare-spike}
Fix \(0<\delta<1/2\).  For every \(\theta>0\) and every \(\rho>0\), there are
independent centered real random variables \(\xi_2,\xi_3\) with
\[
\norm{\xi_2}_{L^2}=\norm{\xi_3}_{L^2}=1,\qquad
\delta\le \norm{\xi_2}_{L^1},\norm{\xi_3}_{L^1}\le 1-\delta,
\]
such that, for the orthogonal unit vectors \(a=e_2\) and \(b=e_3\),
\[
\Prob\left(\bigl||X_a|-|X_b|\bigr|\ge\theta\right)<\rho .
\]
\end{proposition}

\begin{proof}
Set \(r=\delta\).  Choose \(M>r\), to be taken large, and define
\[
        p_M:=\frac{1-r^2}{M^2-r^2}.
\]
Let \(\xi_M\) be symmetric and satisfy
\[
\Prob\{\xi_M=r\}=\Prob\{\xi_M=-r\}=\frac{1-p_M}{2},
\qquad
\Prob\{\xi_M=M\}=\Prob\{\xi_M=-M\}=\frac{p_M}{2}.
\]
Then \(\E\xi_M=0\), and
\[
\E\xi_M^2=(1-p_M)r^2+p_M M^2
        =r^2+p_M(M^2-r^2)=1.
\]
Moreover,
\[
\E|\xi_M|
=(1-p_M)r+p_M M
=r+p_M(M-r)
=r+\frac{1-r^2}{M+r}.
\]
Since \(r=\delta<1/2\), for all sufficiently large \(M\),
\[
        \delta\le \E|\xi_M|\le 1-\delta .
\]
Let \(\xi_2,\xi_3\) be independent copies of \(\xi_M\).  For \(a=e_2\) and
\(b=e_3\), one has \(X_a=\xi_2\) and \(X_b=\xi_3\).  The absolute values of
\(\xi_2\) and \(\xi_3\) take only the values \(r\) and \(M\).  If
\(M-r>\theta\), then
\[
\Prob\left(\bigl||\xi_2|-|\xi_3|\bigr|\ge\theta\right)
=2p_M(1-p_M)\le 2p_M.
\]
As \(M\to\infty\), \(p_M\to0\).  Taking \(M\) large enough to ensure both
\(M-r>\theta\) and \(2p_M<\rho\) proves the claim.
\end{proof}

\begin{remark}[Endpoint \(\delta=1/2\)]
The endpoint is not materially different.  Take \(r<1/2\) and choose
\[
        M=\frac{1-r^2}{1/2-r}-r.
\]
Then the same construction gives \(\E|\xi_M|=1/2\), while \(M\to\infty\) as
\(r\uparrow1/2\).  Hence the probability of a modulus discrepancy can again be
made arbitrarily small.
\end{remark}

\begin{remark}[Where the compactness proof breaks]\label{rem:break}
The obstruction is exactly the distinction between convergence in probability
and the \(L^1\) convergence of the quadratic observable.  In the construction
above, two independent copies satisfy
\[
\bigl||\xi_2|-|\xi_3|\bigr|\longrightarrow0
\qquad\text{in probability}
\]
as \(M\to\infty\), because the mismatch event has probability
\(2p_M(1-p_M)\to0\).  However,
\[
\E\bigl(|\,|\xi_2|^2-|\xi_3|^2\,|\bigr)
=2p_M(1-p_M)(M^2-r^2)
\longrightarrow 2(1-r^2).
\]
Thus the quadratic probes used in the stable \(L^2\) argument see the rare
high-energy events, while weak convergence and convergence in probability do
not.  A limit-level ``shifted head identification'' cannot recover these
missing quadratic contributions without an additional uniform-integrability or
small-ball hypothesis.
\end{remark}

\begin{corollary}[The finite empirical theorem is false as stated]
\label{cor:finite-false}
For \(0<\delta<1/2\), there are no constants \(C,c,\gamma>0\), depending only
on \(\delta\) and on the law of \(\eta\), for which the following statement
holds for every finite system satisfying the hypotheses in Section~1:
if \(m\ge Cn\), then with probability at least \(1-e^{-cm}\),
\[
\frac1m\sum_{j=1}^m
\left(
  |\langle Z_j,u\rangle|-|\langle Z_j,v\rangle|
\right)^2
\ge
\gamma\,\distpm{u}{v}^{\,2}
\qquad(u,v\in\Sph^{n-1}).
\]
\end{corollary}

\begin{proof}
Suppose such constants existed.  Take \(n=3\) and let \(\xi_2,\xi_3\) be the
rare-spike variables from Proposition~\ref{prop:rare-spike}, with \(M\) still
free.  Put \(u=e_2\) and \(v=e_3\).  Then \(\distpm{u}{v}^2=2\).

Let \(m=\lceil 3C\rceil\), and let
\[
q_M:=\Prob\{|\xi_2|\ne|\xi_3|\}=2p_M(1-p_M).
\]
On the event that no one of the \(m\) samples has
\(|\xi_{2,j}|\ne|\xi_{3,j}|\), the empirical average is zero, hence the desired
lower bound fails.  The probability of this event is
\[
        (1-q_M)^m.
\]
Since \(q_M\to0\), we may choose \(M\) so large that
\[
        (1-q_M)^m>e^{-cm}.
\]
Then the failure probability is larger than the allowed \(e^{-cm}\), a
contradiction.
\end{proof}

\section{A Conditional VC Consequence}

The VC part of the draft is correct once a genuine orthogonal small-ball input
is assumed.  The finite-dimensional reduction should be written pointwise; no
empirical covariance concentration is needed, and such concentration is not
available from the hypotheses in Section~1 alone.

\begin{assumption}[Orthogonal small-ball input]\label{ass:orth-small-ball}
For a random vector \(Z\in\R^n\), assume that there are constants
\(\theta_0,p_0>0\) such that
\[
\inf_{\substack{p,q\in\Sph^{n-1}\\ p\perp q}}
\Prob\left(
  \bigl||\langle Z,p\rangle|-|\langle Z,q\rangle|\bigr|\ge\theta_0
\right)\ge p_0 .
\]
\end{assumption}

\begin{lemma}[VC dimension of the orthogonal gap events]\label{lem:gap-vc}
For fixed \(\theta_0>0\), define
\[
\mathcal E(p,q)
=
\left\{
x\in\R^n:
\bigl||\langle x,p\rangle|-|\langle x,q\rangle|\bigr|^2
\ge \theta_0^2
\right\}.
\]
The class
\[
        \{\mathbf 1_{\mathcal E(p,q)}:p,q\in\Sph^{n-1}\}
\]
has VC-dimension \(O(n)\).
\end{lemma}

\begin{proof}
Write \(A=\langle x,p\rangle\) and \(B=\langle x,q\rangle\).  The condition
\((|A|-|B|)^2\ge\theta_0^2\) is equivalent to the conjunction
\[
        A^2+B^2\ge\theta_0^2
\]
and
\[
        (A^2+B^2-\theta_0^2)^2\ge 4A^2B^2.
\]
These are polynomial threshold conditions of degrees \(2\) and \(4\) in
\((x,p,q)\in\R^n\times\R^n\times\R^n\).  Lemma~\ref{lem:poly-vc} and
Proposition~\ref{prop:vc-intersection} give the claimed \(O(n)\) bound.
\end{proof}

\begin{lemma}[Pointwise reduction from arbitrary pairs to orthogonal pairs]
\label{lem:pointwise-reduction}
Let \(u,v\in\Sph^{n-1}\), and let \(d=\distpm{u}{v}\).  For every
\(x\in\R^n\), there are orthogonal unit vectors \(p,q\in\Sph^{n-1}\), depending
only on \(u,v\), such that
\[
\left(
  |\langle x,u\rangle|-|\langle x,v\rangle|
\right)^2
\ge
\frac{\theta_0^2 d^2}{2}\,
\mathbf 1_{\mathcal E(p,q)}(x).
\]
\end{lemma}

\begin{proof}
The assertion is trivial if \(d=0\).  Since replacing \(v\) by \(-v\) does not
change either \(|\langle x,v\rangle|\) or \(\distpm{u}{v}\), we may assume
\[
        d=\norm{u-v}_2\le\norm{u+v}_2 .
\]
Set
\[
        h=\frac{u+v}{2},\qquad k=\frac{u-v}{2},
\qquad
        \alpha=\norm{h}_2,\qquad \beta=\norm{k}_2=\frac d2 .
\]
Then \(h\perp k\) and \(\alpha\ge\beta\).  If \(\beta=0\), again there is
nothing to prove.  Let
\[
        e=\frac{h}{\alpha},\qquad f=\frac{k}{\beta},
\qquad
        p=\frac{e+f}{\sqrt2},\qquad q=\frac{e-f}{\sqrt2}.
\]
Then \(p,q\) are orthogonal unit vectors.  Put
\[
        A=\langle x,e\rangle,\qquad B=\langle x,f\rangle.
\]
For real numbers \(s,t\),
\[
        \bigl||s+t|-|s-t|\bigr|=2\min(|s|,|t|).
\]
Therefore
\[
\bigl||\langle x,p\rangle|-|\langle x,q\rangle|\bigr|
=\sqrt2\min(|A|,|B|).
\]
On the event \(x\in\mathcal E(p,q)\), this gives
\[
        \min(|A|,|B|)\ge\frac{\theta_0}{\sqrt2}.
\]
Using \(u=h+k=\alpha e+\beta f\) and \(v=h-k=\alpha e-\beta f\), the same
scalar identity yields
\[
\bigl||\langle x,u\rangle|-|\langle x,v\rangle|\bigr|
=2\min(\alpha|A|,\beta|B|)
\ge 2\beta\min(|A|,|B|)
\ge \frac{\theta_0 d}{\sqrt2}.
\]
Squaring proves the lemma.
\end{proof}

\begin{theorem}[Conditional finite-dimensional empirical consequence]
\label{thm:conditional-vc}
Let \(Z\in\R^n\) satisfy Assumption~\ref{ass:orth-small-ball}, and let
\(Z_1,\ldots,Z_m\) be i.i.d. copies of \(Z\).  There are constants
\(C_1,c_1>0\), depending only on \(p_0\) and on the absolute constants in the
VC estimates, such that if \(m\ge C_1n\), then with probability at least
\(1-e^{-c_1m}\),
\[
\frac1m\sum_{j=1}^m
\left(
  |\langle Z_j,u\rangle|-|\langle Z_j,v\rangle|
\right)^2
\ge
\frac{\theta_0^2p_0}{4}\,\distpm{u}{v}^{\,2}
\]
for every \(u,v\in\Sph^{n-1}\).
\end{theorem}

\begin{proof}
By Lemma~\ref{lem:gap-vc}, the class
\[
\mathcal F=\{\mathbf 1_{\mathcal E(p,q)}:p,q\in\Sph^{n-1}\}
\]
has VC-dimension at most \(K n\), where \(K\) is absolute.  Choose \(C_1\)
large and \(c_1\) small enough that, whenever \(m\ge C_1n\),
\[
C\sqrt{\frac{Kn}{m}}+\sqrt{\frac{c_1}{2}}\le \frac{p_0}{2},
\]
where \(C\) is the constant in Proposition~\ref{prop:vc-lln}.  Applying
Proposition~\ref{prop:vc-lln} with \(t=c_1m\), we get, with probability at
least \(1-e^{-c_1m}\),
\[
\sup_{p,q\in\Sph^{n-1}}
\left|
\frac1m\sum_{j=1}^m\mathbf 1_{\mathcal E(p,q)}(Z_j)
-
\Prob(Z\in\mathcal E(p,q))
\right|
\le \frac{p_0}{2}.
\]
On this event, Assumption~\ref{ass:orth-small-ball} implies that for every
orthogonal unit pair \(p,q\),
\[
\frac1m\sum_{j=1}^m\mathbf 1_{\mathcal E(p,q)}(Z_j)\ge \frac{p_0}{2}.
\]

Now fix arbitrary \(u,v\in\Sph^{n-1}\).  Let \(p,q\) be the orthogonal unit
pair supplied by Lemma~\ref{lem:pointwise-reduction}.  Summing the pointwise
inequality from that lemma over the samples gives
\[
\frac1m\sum_{j=1}^m
\left(
  |\langle Z_j,u\rangle|-|\langle Z_j,v\rangle|
\right)^2
\ge
\frac{\theta_0^2\distpm{u}{v}^{\,2}}{2}
\cdot
\frac1m\sum_{j=1}^m\mathbf 1_{\mathcal E(p,q)}(Z_j).
\]
The empirical proportion on the right is at least \(p_0/2\), so the desired
bound follows.
\end{proof}

\section{Conclusion}

The VC discretization is sound as a conditional argument from a uniform
orthogonal small-ball estimate.  The missing point is that the one-exception
\(L^1\)-\(L^2\) hypotheses do not provide such a small-ball estimate.  They
give enough information for stable \(L^2\) phase retrieval, where rare
high-energy events are visible to quadratic observables, but they do not
control the probability with which the modulus gap is visible in finitely many
samples.  To recover the originally intended empirical theorem, one needs an
additional assumption excluding the rare-spike obstruction, such as a genuine
Paley--Zygmund input, a uniform \(L_{2+\alpha}\)-to-\(L_2\) condition, or some
other small-ball hypothesis strong enough to imply
Assumption~\ref{ass:orth-small-ball}.

\end{document}
